\documentclass{xStandalone}

\newcommand{\lne}[1]
{
    \draw[-,ultra thin,red] (#1,-0.1)--(#1,0.1);
}

\newcommand{\pnt}[1]
{
    \draw[-,ultra thin,red,fill=red!70!white] (#1,0) circle (0.04);
}

\begin{document}
\begin{tikzpicture}

\draw[gray,ultra thin] (-2.9,-4.1) rectangle (2.9,4.1);

\draw[black!30!white,ultra thin,dashed] (0,-4)--(0,+4);

\draw[gray,very thick] (-2.8,0)--(+2.8,0);

\begin{scope}[decoration={markings,
    mark=
    between positions 0.15 and 0.9 step 0.8cm
    with
    {
        % \lne{-0.3}
        % \pnt{-0.1}
        % \lne{+0.1}
        % \pnt{+0.3}
    }}]
    \draw[-latex,thin,gray] 
    (120:4)-- node[left,pos=0.4] {\tiny 自然光}
    (0,0);
\end{scope}

\begin{scope}[decoration={markings,
    mark=
    between positions 0.2 and 0.95 step 0.8cm
    with
    {
        \lne{-0.3}
        \pnt{-0.1}
        \pnt{+0.1}
        \pnt{+0.3}
    }}]
    \draw[-latex,thick,postaction={decorate}] 
    (0,0)-- node[right,pos=0.15] {\tiny 反射光垂直纸面部分偏振}
    (60:4);
\end{scope}

\begin{scope}[decoration={markings,
    mark=
    between positions 0.2 and 0.95 step 0.8cm
    with
    {
        \lne{-0.3}
        \lne{-0.1}
        \lne{+0.1}
        \pnt{+0.3}
    }}]
    \draw[-latex,thick,postaction={decorate}] 
    (0,0)-- node[right=0.07cm,pos=0.15] {\tiny 折射光平行纸面部分偏振}
    (-70.53:4);
\end{scope}

\path (0,0) node[point]{};


\end{tikzpicture}    
\end{document}

